\documentclass[../../hsm_tok_q.tex]{subfiles}

\begin{document}
    \setcounter{figure}{0}
    \textsf{図\ref{fig/hsm_tok_2016_1st_04_01_q}}で，
    四角形ABCDは平行四辺形であり，
    $\mathrm{AB}=10\,\mathrm{cm}$，$\mathrm{BC}=8\,\mathrm{cm}$，
    $\angle\mathrm{ABC}=60^\circ$である．
    
    辺BC上に点Eをとり，$\mathrm{BE}=3\,\mathrm{cm}$とする．
    
    2つの対角線の交点をOとし，線分AEと線分BDとの交点をFとする．
    
    次の各問に答えよ．
    %
    \vspace{.5\baselineskip}
    {\par \centering
        \setlength{\abovecaptionskip}{5pt}
        \includegraphics[width=.75\linewidth]%
            {\subfix{fig_hsm_tok_2016_1st_04/fig_hsm_tok_2016_1st_04_01_q.pdf}}
            \captionof{figure}{} \label{fig/hsm_tok_2016_1st_04_01_q}
    \par}
    %
    \begin{enumerate}
        \item 線分BDの長さを求めよ．

        \item $\mathrm{△ABE}\equiv\mathrm{△CDF}$であることを証明せよ．

        \newpage

        \item \textsf{図\ref{fig/hsm_tok_2016_1st_04_02_q}}は，
            \textsf{図\ref{fig/hsm_tok_2016_1st_04_01_q}}において，
            線分CFを延長した直線と辺ADとの交点をGとした場合を表している．
            
            $\mathrm{△AEG}$の面積を求めよ．
            %
            \vspace{.5\baselineskip}
            {\par \centering
                \setlength{\abovecaptionskip}{5pt}
                \includegraphics[width=.75\linewidth]%
                    {\subfix{fig_hsm_tok_2016_1st_04/fig_hsm_tok_2016_1st_04_02_q.pdf}}
                    \captionof{figure}{} \label{fig/hsm_tok_2016_1st_04_02_q}
            \par}
        \end{enumerate}

\end{document}
